% Junior Level 1, Week 9 -- answer key
\answersweek{9}{Project Studio}{Answers for workbook chapter 9 and play-at-home card 9.}

\section{Lesson questions}
\begin{answer}[Exit question]
Click the sprite or press the key named in the hat block of that script, for example \enquote{click the balloon} for \enquote{when this sprite clicked} under Balloon1, or \enquote{press space} for \enquote{when space key pressed}.
\end{answer}

\begin{answer}[Reflection]
Any number of colored boxes is fine. Three boxes means the card works (milestone M2). Fewer boxes: plan extra studio time (buffer B2 or the start of week 10). Any honest smiley is right.
\end{answer}

\begin{answer}[Ask before you show (safety talk)]
Priya asks Kai first: \enquote{May I show your card to my sister?} If Kai says no, she does not show it.
\end{answer}

\section{Workbook exercises}
\begin{answer}[Exercise 9.1: My Storyboard (milestone M1)]
Every storyboard is different. Look for three boxes with one idea each: box 1 starts with the green flag (someone says hello and there is a sound); box 2 has something that happens again and again (a dance or animation: a repeat); box 3 starts with a click or a key and has a surprise (a second sprite or a new backdrop). Stick figures are perfect. Too many ideas in one box: move the extras to the back of the sheet.
\end{answer}

\begin{answer}[Exercise 9.2: Build My Card (milestone M2)]
\textbf{What you should see in Mei's sample:} after the green flag, the cat says \enquote{Happy birthday!} for 2 seconds, meows, then changes costume four times, one second apart (a little dance). Clicking the balloon plays a pop and the balloon says \enquote{Surprise!}.

\textbf{What to look for in every card:} \tickbox{} the green flag starts it \quad \tickbox{} at least one repeat \quad \tickbox{} one click or key event \quad \tickbox{} a sound \quad \tickbox{} two sprites or a backdrop change \quad \tickbox{} each box was tested before the next one was started. Full reference: the exercise folder \texttt{exercises/week09/ex-build-my-card/solution.md} and the project teacher guide.
\end{answer}

\begin{answer}[Exercise 9.3: Partner Check]
All six rows say \enquote{yes} for a finished card. A \enquote{not yet} is a good find, not a failure: the owner fixes it next. Common fixes: put a hat block on top of a script that never starts; move a click script to the sprite the storyboard names; put the dance inside a repeat; add a sound block.
\end{answer}

\section{Play-at-home card}
\begin{answer}[Practice the show]
Look for three short sentences in order (start, middle, end), who the card is for (a nickname or \enquote{Grandpa}, never a full name), and one script explained as \enquote{When you click the balloon, it pops and says Surprise}.
\end{answer}

\begin{answer}[Finish a box]
Any finished drawing of the missing box is right. If it was built at home, it should match the drawing.
\end{answer}
